%! Rates of Change in Linear and Quadratic Functions — Unit 1, Lesson 1.3 — Exit Ticket (Answer Key)
\documentclass[10pt]{article}
\usepackage{precalculus-article}
\usepackage{precalculus-key}   % pulls in -boxes; do NOT also load -boxes
\newcommand{\TallMath}[1]{$\displaystyle #1\rule[-1.4em]{0pt}{3.2em}$}

\begin{document}

\pageheader{Unit 1, Lesson 1.3}{Exit Ticket --- Answer Key}
\namedateperiod

\vspace{0.15in}

\begin{enumerate}[leftmargin=1.5em, itemsep=12pt]

\item A function $h(x)$ has the following values.

\smallskip
\begin{tabular}{|c|c|c|c|c|c|}
\hline
$x$    & 0 & 1 & 2 & 3  & 4  \\
\hline
$h(x)$ & 1 & 3 & 7 & 13 & 21 \\
\hline
\end{tabular}

\smallskip
\begin{enumerate}[label=(\alph*), itemsep=6pt]
  \item Compute the AROC (first difference) for each consecutive unit interval.

  \smallskip
  \begin{tabular}{|l|c|}
  \hline
  \textbf{Interval} & \textbf{AROC} \\
  \hline
  $[0, 1]$ & \ans{2} \\
  $[1, 2]$ & \ans{4} \\
  $[2, 3]$ & \ans{6} \\
  $[3, 4]$ & \ans{8} \\
  \hline
  \end{tabular}

  \item Are the second differences constant? \ans{Yes, all equal 2.} \quad What type of function does this suggest? \ans{Quadratic}
\end{enumerate}

\item For a linear function, the AROC from $x = 2$ to $x = 5$ is 4.
What is the AROC from $x = 5$ to $x = 10$? Explain your answer in one sentence.

\vspace{0.1in}
AROC on $[5, 10]$: \ans{4}

\ansline{For a linear function the AROC is constant over any interval, so it must still equal 4.}

\end{enumerate}

\end{document}
